% Ready-to-integrate supplement to the 2026-10-06 leading-coefficient audit.
% No changes to the canonical BFSS v2.1 manuscript are made by this file.
% Requires amsmath and amssymb. All labels have the auditadd: prefix.

\section{Explicit slow Clifford contraction and finite sector transport}
\label{auditadd:section}
This supplement expands two compressed steps in Sections 3 and 6.1 of
\emph{BFSS leading coefficient cancellation audit}, 6 October 2026.
It concerns the uncut leading coefficients at fixed finite $N$, using the
kinetic-one, reduced-kinetic, and frozen-pair definitions of Appendices
A, B, and D of the supplied v2.1 manuscript. The identities below do not
establish an all-vector reserve, form-domain compatibility, the global
Hardy theorem, or constants uniform in $N$.

\subsection{The pair contraction before any slow Clifford trace}
Fix an oriented edge with block sizes $n,m$, put $d=nm$, and choose its
spatial frame so that the center difference is $r e_9$, $r>0$.
Let $t_c$ be a trace-orthonormal real basis of the block-diagonal slow
color space, including the normalized relative center. Its Hermitian
action on $\operatorname{Hom}(\mathbb C^m,\mathbb C^n)$ is
\[
 T_c q=(t_c)_{aa}q-q(t_c)_{bb},\qquad
 K_{cc'}=\operatorname{tr}_{\mathbb C^d}(T_cT_{c'}).
\]
Other blocks can be included in an ambient orthonormal slow basis: their
restriction merely adds zero actions and an orthogonal change of the
center basis. In particular the following calculation applies to every
partition. Since $T_c$ and $T_{c'}$ are Hermitian, $K$ is real symmetric.

Let $a,h,k$ label the $2d$ real off-block color generators. With
$[t_A,t_B]=i f_{ABC}t_C$, set
\[
 (F_c)_{ah}=f_{ahc}.
\]
The symbol $F_c$ in this subsection denotes a color-action matrix, not
the inverse-Gauss coefficient $F$ in Appendix D. On the underlying real
space of $\mathbb C^d$, $F_c$ is the realification of $iT_c$. Let $J$ be
multiplication by $i$ on that real space. Then
\begin{align}
 \sum_{a,h}(F_c)_{ah}(F_{c'})_{ah}&=2K_{cc'},
       \tag{A0a}\label{auditadd:color-real}\\
 \sum_{a,h,k}(F_c)_{ah}(F_{c'})_{ak}J_{hk}&=0.
       \tag{A0b}\label{auditadd:color-polarization}
\end{align}
Indeed the first expression is $\operatorname{tr}_{\mathbb R}(F_c^T F_{c'})$.
For the second, $F_c^T F_{c'}$ is the realification of $T_cT_{c'}$;
contracting with $J$ gives, up to the immaterial orientation sign,
$2\operatorname{Im}\operatorname{tr}_{\mathbb C}(T_cT_{c'})=0$.
This vanishes for each $c,c'$, not just after summing $c=c'$.

Write $\xi_{h\alpha}$ for fast and $\eta_{c\beta}$ for slow Majoranas,
with $\{\xi_i,\xi_j\}=\delta_{ij}$,
$\{\eta_i,\eta_j\}=\delta_{ij}$, and $\{\xi_i,\eta_j\}=0$.
The real symmetric Spin(9) matrices satisfy
$\{\Gamma_i,\Gamma_j\}=2\delta_{ij}I_{16}$.
In the stated orientation the centered covariance is
\[
 \langle\xi_{h\alpha}\xi_{k\gamma}\rangle_f
 =\frac12\bigl(\delta_{hk}\delta_{\alpha\gamma}
              -iJ_{hk}(\Gamma_9)_{\alpha\gamma}\bigr).
 \tag{A1}\label{auditadd:covariance}
\]
Changing the orientation changes both occurrences of the complex
structure consistently and does not affect the results.

The mixed slow--fast part of the off-block Gauss spin is
\[
 G_a=-i\sum_{h,c,\beta}(F_c)_{ah}\xi_{h\beta}\eta_{c\beta}.
\]
Move every slow generator past the second fast generator before taking
the fast expectation. The minus sign from this move cancels
$(-i)^2$. Equations \eqref{auditadd:color-real}--\eqref{auditadd:covariance}
give the following equality in the \emph{untraced} slow Clifford algebra:
\begin{align}
 \sum_a\langle G_a^2\rangle_f
 &=\sum_{c,c',\beta,\gamma}
       Q^G_{c\beta,c'\gamma}\eta_{c\beta}\eta_{c'\gamma},
 &Q^G_{c\beta,c'\gamma}&=K_{cc'}\delta_{\beta\gamma}.
 \tag{A2}\label{auditadd:gauss-coefficient}
\end{align}
There is no residual polarized bilinear: its color coefficient is
exactly \eqref{auditadd:color-polarization}. Applying CAR to the real
symmetric coefficient in (A2) gives
\[
 \sum_a\langle G_a^2\rangle_f
 =8\sum_cK_{cc}I=8d(n+m)I.
 \tag{A3}\label{auditadd:gauss-scalar}
\]
The last equality follows from the trace-normalized endpoint Casimirs:
\[
 \sum_{c\,\mathrm{internal}}K_{cc}=(d-1)(n+m),\qquad
 K_{\mathrm{relative\ center},\mathrm{relative\ center}}=n+m.
\]
The centered inverse-Gauss square multiplies (A3) by $r^{-2}$.

For the source, let $y_{ia}$, $1\le i\le8$, be the real transverse
oscillator coordinates; their covariance is
$\langle y_{ia}y_{jb}\rangle=\delta_{ij}\delta_{ab}/(2r)$.
In the conventions of the audit the mixed Yukawa source is
\[
 S_e(0)=\left[-2i\sum_{i=1}^8\sum_{a,h,c,\alpha,\beta}
 y_{ia}(\Gamma_i)_{\alpha\beta}(F_c)_{ah}
 \xi_{h\alpha}\eta_{c\beta}\right]U_e(0).
 \tag{A4}\label{auditadd:source}
\]
Here $U_e(0)$ is the normalized bosonic/Fock vacuum embedding; (A4)
is a map from the slow module, so $S_e(0)^*S_e(0)$ leaves that module
untraced. Moving slow past fast generators as above and contracting
$y,\xi$ gives
\begin{align}
 S_e(0)^*S_e(0)
 &=\sum_{c,c',\beta,\gamma}
    Q^Y_{c\beta,c'\gamma}\eta_{c\beta}\eta_{c'\gamma},
 & Q^Y_{c\beta,c'\gamma}
 &=\frac{16}{r}K_{cc'}\delta_{\beta\gamma}.
 \tag{A5}\label{auditadd:yukawa-coefficient}
\end{align}
For clarity, the unpolarized coefficient is
$4(2r)^{-1}\cdot\tfrac12\cdot2K_{cc'}$
times $\sum_{i=1}^8\Gamma_i^T\Gamma_i=8I$.
The polarized coefficient contains
$\sum_{i=1}^8\Gamma_i^T\Gamma_9\Gamma_i=-8\Gamma_9$
but vanishes by \eqref{auditadd:color-polarization}.
Consequently CAR yields the exact operator identity
\[
 S_e(0)^*S_e(0)=\frac{128d(n+m)}r I.
 \tag{A6}\label{auditadd:source-scalar}
\]
The source has one bosonic and one fermionic excitation on this edge,
each of energy $2r$. Thus its centered resolvent contraction equals
$32d(n+m)r^{-2}I$, as claimed in Section 3 of the audit. This calculation
establishes the slow-Clifford reduction directly; a trace over the slow
module would not suffice to establish it.

\subsection{Transport preserving actual excitation sectors}
Fix the centers and use $A(t)=tA$, $0\le t\le1$. On each edge, let the
real positive kinetic and potential matrices be $G_e(t)$ and $V_e(t)$.
Set
\[
 \Omega_e=G_e^{-1/2}(G_e^{1/2}V_eG_e^{1/2})^{1/2}G_e^{-1/2},
 \qquad \Omega_eG_e\Omega_e=V_e.
\]
All these expressions concern finite coefficient matrices. Define the
unitary map from the fixed dimensionless oscillator space to the actual
bosonic fiber by
\[
 (\mathcal U_e\psi)(Y_e)
   =(\det\Omega_e)^{1/4}\psi(\Omega_e^{1/2}Y_e).
\]
Writing $x_e=\Omega_e^{1/2}Y_e$ and
$M_e=\Omega_e^{1/2}G_e\Omega_e^{1/2}$ gives exactly
\begin{align}
 \mathcal U_e^*H_{B,e}\mathcal U_e
 &=p_{x_e}^TM_ep_{x_e}+x_e^TM_ex_e,\nonumber\\
 \mathcal U_e^*(H_{B,e}-E_{B,e})\mathcal U_e
 &=2\sum_{j,k}a_j^*(M_e)_{jk}a_k,
 &E_{B,e}&=\operatorname{tr}_{\mathbb R}M_e,
 \tag{A7}\label{auditadd:boson-transport}
\end{align}
where $a_j=(x_j+\partial_{x_j})/\sqrt2$. Thus total bosonic excitation
number is invariant even when $G_e,V_e$ do not commute. There is no
choice of eigenvectors or division by differences of internal
frequencies. At $t=0$, $M_e=r_eI$.

For fermions let $D_e(t)$ be the Hermitian complex bifundamental mass
and put
\[P_e(t)=1_{(-\infty,0)}(D_e(t)).\]
With the external gap retained,
solve the finite-dimensional unitary transport equation
\[
 \dot V_e^{\mathrm F}=[\dot P_e,P_e]V_e^{\mathrm F},\qquad
 V_e^{\mathrm F}(0)=I.
 \tag{A8}\label{auditadd:kato}
\]
Then $P_e(t)=V_e^{\mathrm F}(t)P_e(0)V_e^{\mathrm F}(t)^*$ and
$\widetilde D_e=V_e^{\mathrm F*}D_eV_e^{\mathrm F}$ commutes with
$P_e(0)$. Second quantization transports the occupied determinant to
the fixed reference determinant. Relative to particles in
$1-P_e(0)$ and holes in $P_e(0)$, the positive excitation matrix in the
kinetic-one normalization is
\[
 B_e^{\mathrm F}=2\left(
 \widetilde D_e\big|_{1-P_e(0)}\ \oplus\
 [-\widetilde D_e\big|_{P_e(0)}]^T\right).
 \tag{A9}\label{auditadd:fermion-transport}
\]
The transpose acts on the dual hole space. The shifted fermionic
Hamiltonian is its fermionic second quantization and preserves total
particle-plus-hole excitation number. Its centered value is $2r_eI$
on the one-excitation space.

Uniform pair functional calculus on the normalized thinness ball gives
\[
 \|M_e-r_eI\|\le C_N\kappa r_e,\quad M_e\ge c r_eI,
 \qquad
 \|B_e^{\mathrm F}-2r_eI\|\le C_N\kappa r_e,\quad
 B_e^{\mathrm F}\ge c r_eI.
 \tag{A10}\label{auditadd:pair-matrices}
\]
These bounds use the external positive/negative gap of $D_e$ and the
positive bosonic gap; no internal eigenvalue gap is assumed. The
transport is pairwise. In particular its bound never involves an
unrelated $r_f/r_e$.

On the fixed tensor product, let $N_e$ be total bosonic plus fermionic
excitation number on edge $e$. The actual shifted fast Hamiltonian is
the sum of the transported pair excitation Hamiltonians. Each
spectator edge has zero energy on its transported vacuum because its
own ground energy has been subtracted. Hence the spaces
\[
 \mathcal E_e=\{N_e=2,\ N_h=0\ (h\ne e)\},\qquad
 \mathcal E_t=\{N_e=N_f=N_g=1,\ N_h=0\ (h\notin t)\}
\]
are finite-dimensional, mutually orthogonal, and invariant. The pair
source lies in the smaller one-boson/one-fermion subspace of
$\mathcal E_e$. The monomial inventory in Appendix D(2.6) places each
triangle source in $\mathcal E_t$. Different such monomials can mix
\emph{within} a sector; the argument neither discards that mixing nor
requires individual internal eigenmodes to remain fixed.

Since at most two quanta occur on an edge sector and one on each edge
of a triangle, (A10) proves
\begin{align}
 H_{e,A}&\ge c r_eI,&
 \|H_{e,A}-H_{e,0}\|&\le C_N\kappa r_e,\nonumber\\
 H_{t,A}&\ge c s_tI,&
 \|H_{t,A}-H_{t,0}\|&\le C_N\kappa s_t,
 &s_t&=\sum_{h\in t}r_h.
 \tag{A11}\label{auditadd:sector-bounds}
\end{align}
Here all differences are taken after the specified pair transports.
The resolvent identity now gives $C_N\kappa/r_e$ and
$C_N\kappa/s_t$, respectively, for the inverse differences. Combining
these with the actual transported source difference bounds in Section
6.1 of the audit proves its displayed Schur-complement difference
estimate. In particular the sources have not been defined by
transporting their centered values: one first forms $H_1(A)U(A)$ and
then uses the transports to compare it with the centered source.

\subsection{Derivative scope and reproducibility}
The additional derivatives used for the multiplier in Section 6 of the
audit are derivatives of the finite coefficient matrices and their
spatial/internal jets. An internal derivative brings its incident
$r_e^{-1}$. A covariant center derivative differentiates $r_e$, the
unit direction, and the normalized internal matrix, again with an
incident inverse gap. Along $tA$, each normalized jet changes by
$O(C_N\kappa)$, using one further bounded derivative of the same matrix
functions. This explains why coincident internal eigenvalues cause no
singularity. The finite root-path inventory and its triangle
inequalities remain necessary; smoothness by itself would not control
unrelated gap ratios.

The supplementary \texttt{check\_pair\_clifford.py} constructs the
untraced slow quadratic coefficient matrices in (A2) and (A5) for
$(n,m)=(1,1),(2,1),(2,2),(3,1)$ at $r=1$. It compares every coefficient
with the displayed Gram-matrix formulas and tests the vanishing
antisymmetric coefficients, not just the resulting scalar trace.
It uses floating-point arithmetic and reuses the original test's
Spin(9) and color-basis constructors, with separate contraction code.
These finitely many tests are regression evidence; the arbitrary-block
result is the preceding algebraic calculation.

The original symbolic script establishes exact rational reductions of
its supplied triangle formulas. The original floating script tests four
centered singleton-triangle geometries, reconstructing their triangle
Gaussian/Fock contractions while supplying the singleton pair
Born--Huang and spin terms as known constants. Neither original script
computes the nonzero-$A$ operator, verifies the arbitrary-block
slow-Clifford calculation, or tests the required slow derivatives.
No numerical nonzero-$A$ test is supplied by this supplement either;
(A7)--(A11) are analytic statements under the explicitly stated
frozen-pair hypotheses. Nothing in these checks certifies the
all-vector or global-theorem steps excluded at the start.
